\documentclass[aps,prapplied,reprint,groupedaddress,twocolumn,superscriptaddress,linenumbers]{revtex4-2}
\usepackage{CJK}

\usepackage{graphicx}
\usepackage{dcolumn}
\usepackage{bm}
\usepackage{amsmath}
\usepackage{siunitx}
\usepackage{enumitem}
\usepackage{booktabs}

\begin{document}

\preprint{}

\title{Quantum-referenced refractive-index calibration of air at 1762 nm for environmental phase compensation in quantum optical interfaces}

\author{Wei Wu}
\affiliation{Albert-Ludwigs-Universität Freiburg, Physikalisches Institut, Hermann-Herder-Straße 3, 79104 Freiburg, Germany}
\affiliation{EUCOR Centre for Quantum Science and Quantum Computing, University of Freiburg, Hermann-Herder-Straße 3, 79104 Freiburg, Germany}
\author{Tobias Sch{\"a}tz}
\affiliation{Albert-Ludwigs-Universität Freiburg, Physikalisches Institut, Hermann-Herder-Straße 3, 79104 Freiburg, Germany}
\affiliation{EUCOR Centre for Quantum Science and Quantum Computing, University of Freiburg, Hermann-Herder-Straße 3, 79104 Freiburg, Germany}
\author{Ulrich Warring}
\affiliation{Albert-Ludwigs-Universität Freiburg, Physikalisches Institut, Hermann-Herder-Straße 3, 79104 Freiburg, Germany}

\date{\today}

\begin{abstract}
Free-space optical paths accumulate phase errors from environmentally driven refractive-index fluctuations. Models of the atmospheric refractivity are used to predict and compensate these fluctuations, but their accuracy at a specific wavelength is generally not verified by direct measurement. This work presents a measurement of the refractive index of air at \SI{1762}{\nano\meter}, the clock transition of $^{138}\mathrm{Ba}^+$. The measurement is referenced to a single trapped ion and is based on \num{145784} synchronized measurements collected over eight months with a dual-wavelength Michelson interferometer and co-located environmental sensors. The data yield the coefficients $\partial n/\partial X$ for temperature, relative humidity, and pressure. Compared with the Mathar model evaluated on the identical observed conditions, the measured temperature and pressure coefficients agree to within \SI{0.35}{\percent} and \SI{0.93}{\percent}; a Kramers-Kronig analysis of the water-vapor spectrum does not predict an enhancement over the broadband model at this wavelength. The measured humidity coefficient is smaller in magnitude than the same-condition Mathar surrogate by \SI{24.8}{\percent}. This humidity difference is statistically significant in the full dataset, but not within the model validity range; its direction is opposite to the sign expected from water-vapor dispersion, and its magnitude is nearly spanned by the gain uncertainty of the humidity sensor. It is therefore reported as an unresolved systematic. The dominant discrepancy between the measured refractive index and the raw Mathar prediction is a static offset of \SI{4.16}{\radian} per meter; after this offset is removed, the residuals of the empirical model and the Mathar model agree within approximately \SI{5}{\percent}. The dataset of synchronized measurements is made openly available for independent reanalysis.
\end{abstract}

\maketitle

\section{Introduction}
\label{sec:intro}

Optical interconnects between remote trapped-ion quantum nodes require phase-stable links for protocols that depend on the first-order coherence of the optical fields, such as free-space optical clock comparisons \cite{giorgetta2013optical,shen2022free} and coherent detection of transmitted fields. Photon-counting protocols based on post-selection, such as atomic-ensemble quantum memories \cite{duan2001long} and Barrett-Kok-type schemes \cite{barrett2005efficient}, are insensitive to the common optical phase accumulated along the link and therefore do not require phase-stable optical paths. For $^{138}\mathrm{Ba}^+$ based quantum networks requiring phase-stable links, the operating wavelength is \SI{1762}{\nano\meter}, where the $6S_{1/2} \rightarrow 5D_{5/2}$ electric quadrupole transition provides the physical two-level system \cite{dietrich2009barium,arnold2020precision}. The refractive index of air along the optical path depends on temperature ($T$), humidity ($H$), and pressure ($P$), and its fluctuations translate directly into phase errors of the transmitted field.

Compensation of these phase errors requires the refractive index's environmental dependence at the operating wavelength. Atmospheric refractivity is described by broadband formulations, most commonly the Edl\'{e}n, Ciddor, and Mathar models \cite{edlen1966refractive,ciddor1996refractive,mathar2007refractive}. These models are fitted to refractivity data over wide spectral intervals and interpolate between sparse calibration wavelengths via dispersion formulas; their accuracy at a single wavelength is therefore not guaranteed a priori. The Edl\'{e}n and Ciddor formulations are calibrated over the visible and near-infrared (\SIrange{0.4}{1.3}{\micro\meter}), with stated validity extending only to about \SI{1.7}{\micro\meter}, leaving \SI{1762}{\nano\meter} outside their calibrated range. The Mathar model extends the dispersion formalism into the mid-infrared and is the model applicable at \SI{1762}{\nano\meter}. Experimental validation of refractivity models at this wavelength has, to our knowledge, not been reported.

Here we report a quantum-referenced measurement of the refractive index of air at \SI{1762}{\nano\meter} over an eight-month campaign, yielding the environmental coefficients $\partial n/\partial X$ ($X = T, H, P$). We compare the measured coefficients with the Mathar model evaluated on the identical observed conditions. The temperature and pressure coefficients agree with the model to within \SI{0.35}{\percent} and \SI{0.93}{\percent}. The humidity coefficient is smaller in magnitude than the same-condition Mathar surrogate by \SI{24.8}{\percent}; the difference is statistically significant in the full dataset but not within the model validity range, and its interpretation is discussed below. We also report the residual analysis of the empirical model relative to the Mathar formulation and the associated uncertainty budget.

\begin{figure}[htbp]
\centering
\includegraphics[width=0.45\textwidth]{figures/figure1_3d.png}
\caption{Schematic of the experimental setup for quantum-referenced precision refractometry.
\textbf{(Blue)} Measurement node: The \SI{1762}{\nano\meter} probe laser is stabilized to an ultra-low-expansion (ULE) optical cavity for short-term stability \cite{yum2017optical}, and then referenced to the electric quadrupole transition ($6S_{1/2} \rightarrow 5D_{5/2}$) of a single trapped $^{138}\mathrm{Ba}^+$ ion, which sets the optical frequency of this channel. The \SI{780}{\nano\meter} reference laser is stabilized via Doppler-free saturated absorption spectroscopy using a $^{87}$Rb vapor cell. A co-located BME280 sensor monitors temperature ($T$), humidity ($H$), and pressure ($P$) at the interferometer site. A global positioning system (GPS)-disciplined oscillator provides a \SI{10}{\mega\hertz} timebase for synchronization.
\textbf{(Red)} Interferometer node: A dual-wavelength Michelson interferometer is placed in a non-temperature-stabilized building corridor outside the laboratory in Freiburg, Germany, exposing the optical path to ambient variations of temperature, humidity, and pressure.
\textbf{(Yellow)} Database: Correlates the dual-wavelength interferometer data with the co-located environmental sensor readings into synchronized time series, from which the refractive index and its environmental coefficients are derived.
}
\label{fig:architecture}
\end{figure}

\section{Experimental Methods}
\label{sec:methods}

The experimental system is organized into three functional blocks, color-coded in Fig.~\ref{fig:architecture}. The measurement node (blue) provides frequency-referenced laser light and records the environmental conditions at the interferometer site. The \SI{1762}{\nano\meter} probe laser is stabilized to an ultra-low-expansion (ULE) optical cavity (finesse 360,000) and referenced to the electric quadrupole transition ($6S_{1/2} \rightarrow 5D_{5/2}$) of a single trapped $^{138}\mathrm{Ba}^+$ ion, yielding short-term stability $\delta f < \SI{1}{\kilo\hertz}$ \cite{yum2017optical,arnold2020precision}. The \SI{780}{\nano\meter} reference laser is stabilized to the $^{87}$Rb D$_2$ line via Doppler-free saturated absorption spectroscopy ($\delta f < \SI{1}{\mega\hertz}$). A BME280 sensor \cite{bme280} co-located with the interferometer monitors temperature $T$, relative humidity $H$, and pressure $P$ (accuracies $\pm\SI{1.0}{\degreeCelsius}$, $\pm\SI{3}{\percent}$ RH, and $\pm\SI{1}{\hecto\pascal}$). Both lasers share a GPS-disciplined \SI{10}{\mega\hertz} timebase, which stamps every measurement for synchronization; environmental sensors were sampled at \SI{2}{\hertz}.

The interferometer node (red) is a dual-wavelength Michelson interferometer (modified LM10 design of the U.S. National Institute of Standards and Technology (NIST) \cite{fox1999reliable}) measuring over a common round-trip optical path $L = \SI{4.2}{\meter}$ (total propagation). The directly measured quantity is the ratio $R(t)$ of the interference fringe counts accumulated at the two wavelengths, recorded by a frequency counter over 10 to 20 s per sample. Because both beams traverse the same air path, vibrations and drifts common to the two channels cancel in $R(t)$, isolating the differential refractivity of the air. The refractive index at \SI{1762}{\nano\meter} is a derived quantity,

\begin{equation}
n_{1762}(t) = \frac{f_{780}}{f_{1762}}\; \frac{n_{780}\big(T(t),H(t),P(t)\big)}{R(t)},
\label{eqn:wavemeter}
\end{equation}

where $f_{780}$ and $f_{1762}$ are the optical frequencies of the two channels, both fixed by atomic references (the $^{87}$Rb D$_2$ line and the $^{138}\mathrm{Ba}^+$ quadrupole transition, respectively), and $n_{780}$ is obtained from the Ciddor model \cite{ciddor1996refractive} evaluated at the co-located sensor readings.

Equation~(\ref{eqn:wavemeter}) makes explicit what is measured and what is assumed. The directly measured quantity is the fringe-count ratio $R(t)$; the inferred $n_{1762}$ is differential, anchored to the Ciddor model at \SI{780}{\nano\meter} \cite{ciddor1996refractive}, whose input in turn is provided by the co-located BME280 sensor. The inferred $n_{1762}$ therefore inherits any constant offset of the 780~nm anchor as an additive constant; this appears in our data as the static offset discussed in Sec.~\ref{sec:results}. The environmental coefficients $\partial n/\partial X$, which are the quantities of interest here, are unaffected by such a constant because it is absorbed by the fitted intercept. The inferred $n_{1762}$ is the path-averaged refractivity over the \SI{4.2}{\meter} round-trip, whereas $T$, $H$, and $P$ are single-point readings at one end of the path; this spatial mismatch is the dominant contribution to the measurement uncertainty budget (Supplemental Material). The interferometer and the sensor sit in a non-temperature-stabilized building corridor (temperature \SIrange{17.5}{34.6}{\degreeCelsius}, relative humidity \SIrange{18.9}{45.3}{\percent}, pressure \SIrange{964.6}{1006.8}{\hecto\pascal} over the campaign), while the quantum frequency references are operated in a temperature-stabilized laboratory.

Samples recorded while either laser was unlocked were identified by fringe-count ratios inconsistent with locked operation and were removed during preprocessing. This retention criterion, together with the segment structure of the campaign (104 segments between January and August 2025, duty cycle near 15\%, see Supplemental Material), defines the dataset of \num{145784} synchronized measurements \cite{wu_2026_18800166}.

The refractive index at \SI{1762}{\nano\meter} is modeled as a linear function of the environmental variables,

\begin{equation}
n(T,H,P) = n_0 + \alpha_T T + \alpha_H H + \alpha_P P
\end{equation}

where $n_0$ is a fitted intercept and $\alpha_j$ are the environmental coefficients determined by multivariate ordinary least-squares (OLS) regression on the synchronized time series. Statistical robustness is validated through a suite of diagnostics detailed in Supplemental Material, including moving-block bootstrap resampling with a block length of 156 samples (twice the measured $1/e$ residual decorrelation time) and $B = 5000$ replications \cite{kunsch1989jackknife,politis1994stationary}, and heteroskedasticity- and autocorrelation-consistent (HAC) standard errors with a Bartlett kernel matched to the same lag \cite{newey1987hypothesis,andrews1991heteroskedasticity} (agreeing with the bootstrap estimates within 3\%). The predictive performance of the linear surrogate is assessed with a blocked out-of-sample validation: the data are divided into contiguous time blocks assigned to five folds, and within each fold the surrogate is fitted on the training blocks and evaluated on the held-out blocks, with identical folds and no global preprocessing before the split (Supplemental Material).

\begin{figure}[htbp]
\centering
\includegraphics[width=0.45\textwidth]{figures/data_preview.png}
\caption{Synchronous measurement over a 14 h daytime period. \textbf{(a)} Fringe-count ratio $R$ from the dual-wavelength interferometer, the directly measured quantity of Eq.~(\ref{eqn:wavemeter}). \textbf{(b)} Refractive index at \SI{1762}{\nano\meter} derived from (a) via Eq.~(\ref{eqn:wavemeter}), referenced to the Ciddor model at \SI{780}{\nano\meter}. \textbf{(c)} Ambient temperature, \textbf{(d)} relative humidity, and \textbf{(e)} pressure at the interferometer site, recorded by the co-located BME280. The diurnal anti-correlation between (c) and (d) is the physically expected behaviour. Gaps near 16:00 result from laser-unlock events.}
\label{fig:timeseries}
\end{figure}

\section{Results and Validation}
\label{sec:results}

The synchronized dataset of the derived refractive index and the related environmental parameters, shown in Fig.~\ref{fig:timeseries}, spans eight months (January to August 2025) as a series of 104 observation segments with an overall duty cycle near 15\%. The mean temperature, relative humidity, and pressure over the retained observations are \SI{28.9}{\degreeCelsius}, \SI{30.2}{\percent}, and \SI{986.4}{\hecto\pascal}, with standard deviations of \SI{3.4}{\degreeCelsius}, \SI{3.9}{\percent}, and \SI{4.5}{\hecto\pascal}, respectively; time-weighted means, which account for the non-uniform temporal sampling, are \SI{25.6}{\degreeCelsius}, \SI{31.0}{\percent}, and \SI{987.9}{\hecto\pascal}. Details of the temporal coverage are provided in the Supplemental Material.

The refractive index time series, derived from the measured fringe-count ratio via Eq.~(\ref{eqn:wavemeter}), tracks the environmental fluctuations with clear diurnal periodicity. Data gaps, visible in Fig.~\ref{fig:timeseries}, result from ion reloading cycles and occasional interferometer realignment; all valid time points from the full dataset (Section~\ref{sec:methods}) are used for the regression analysis.

Multivariate regression following the linear model defined in Section~\ref{sec:methods} yields the environmental coefficients listed in Table~\ref{tab:coefficients}. To place the comparison with the Mathar model on equal footing, we evaluate the Mathar formulation at the identical $(T,H,P)$ rows of the campaign and fit the same linear surrogate model to both datasets, as detailed in Supplemental Material. The temperature and pressure coefficients agree with this same-condition Mathar surrogate within \SI{0.35}{\percent} and \SI{0.93}{\percent}, respectively. In contrast, the humidity coefficient is smaller in magnitude than the same-condition Mathar surrogate by \SI{24.8}{\percent}. The humidity difference is statistically significant in the full campaign, but a Kramers-Kronig estimate at the mean campaign conditions (Supplemental Material) likewise shows no enhancement over the broadband model, and the humidity-sensor gain systematic nearly spans the observed difference; it is therefore reported as an unresolved systematic pending independent humidity sensor calibration.

\begin{table}[htbp]
\centering
\caption{Measured and predicted environmental coefficients at
         \SI{1762}{\nano\meter}. Mathar surrogate values are obtained
         by fitting the same linear model to the Mathar formulation
         evaluated on the identical observed $(T,H,P)$ rows
         (Supplemental Material). Uncertainties are moving-block
         bootstrap standard errors ($B = 5000$, block length 156
         samples); systematic sensor contributions are provided in
         Supplemental Material.}
\label{tab:coefficients}
\begin{tabular}{lcccc}
\hline
Coefficient & This work & Mathar surrogate & Difference & $u_{\rm stat}$ \\
\hline
$\alpha_T$ (\si{K^{-1}})      & $-8.8474\times 10^{-7}$ & $-8.8164\times 10^{-7}$ & $-0.35\%$ & $1.41\times10^{-9}$ \\
$\alpha_H$ (\%RH$^{-1}$)      & $-1.3152\times 10^{-8}$ & $-1.7486\times 10^{-8}$ & $+24.8\%$ & $1.15\times10^{-9}$ \\
$\alpha_P$ (\si{hPa^{-1}})    & $+2.5949\times 10^{-7}$ & $+2.5710\times 10^{-7}$ & $+0.93\%$ & $9.85\times10^{-10}$ \\
\hline
\end{tabular}
\end{table}

The proximity of \SI{1762}{\nano\meter} to water-vapor absorption features, with the two nearest H$_2$O lines at \SI{1761.0405}{\nano\meter} and \SI{1762.4852}{\nano\meter} \cite{gordon2022hitran2020}, motivates a Kramers-Kronig analysis of the humidity dispersion at this wavelength, detailed in Supplemental Material. The line-by-line prediction from the HITRAN2020 water-vapor spectrum evaluated at the mean campaign conditions yields a humidity coefficient of order $10^{-8}$\,\%RH$^{-1}$ without enhancement over the broadband model at this wavelength. This is consistent with the same-condition comparison above, where the measured humidity coefficient is found to be smaller in magnitude than the Mathar surrogate rather than enhanced. Given the unresolved sensor gain systematic, a quantitative attribution of the residual humidity difference to a wavelength-specific physical effect is not possible with the present data.

The environmental coefficients are converted to phase-shift sensitivity per unit environmental change for $L = \SI{1}{\meter}$ by multiplication with $2\pi L / \lambda$. The resulting engineering quantities are listed in Table~\ref{tab:phase_shift}.

\begin{table}[htbp]
\centering
\caption{Phase-shift sensitivities of free-space optical path per meter
         at \SI{1762}{\nano\meter}, derived from the measured environmental
         coefficients.}
\label{tab:phase_shift}
\begin{tabular}{lcc}
\hline
Environmental parameter & $\partial\phi/\partial X$ & Unit \\
\hline
Temperature          & 3.16  & \si{rad\,m^{-1}\,K^{-1}} \\
Relative humidity    & 0.047 & \si{rad\,m^{-1}\,(\%RH)^{-1}} \\
Pressure             & 0.93  & \si{rad\,m^{-1}\,hPa^{-1}} \\
\hline
\end{tabular}
\end{table}

Table~\ref{tab:phase_shift} lists the resulting phase-shift sensitivities per unit environmental change. The residual analysis of the two models is as follows. The empirical model achieves a residual standard deviation $\sigma_n = 1.84 \times 10^{-7}$, corresponding to a coefficient of determination $R^2 = 0.996$. The Mathar formulation evaluated at the same wavelength produces a root-mean-square residual error of $1.18 \times 10^{-6}$, dominated by a static offset of \num{1.166e-6} (Fig.~\ref{fig:validation}(b)). Since this offset is constant over the campaign, subtracting it leaves a residual standard deviation of \num{1.93e-7}. On this offset-subtracted basis, the empirical model improves the residual phase error by \SI{4.9}{\percent} (amplitude) and \SI{9.5}{\percent} (variance) over the Mathar model. This comparison is summarized in Fig.~\ref{fig:validation}. The in-sample residual of \num{1.84e-7} is representative of the predictive performance: a blocked out-of-sample validation (Supplemental Material) gives a held-out root-mean-square residual of $(1.84 \pm 0.02) \times 10^{-7}$ for the empirical surrogate, confirming that the fit does not overstate its predictive accuracy.

An expanded uncertainty budget is provided in Supplemental Material. The dominant contribution comes from the spatial and temporal mismatch between the environmental sensor and the optical path ($\delta n = 1.84 \times 10^{-7}$). The combined expanded uncertainty with coverage factor $k = 2$ is $3.68 \times 10^{-7}$; the laser frequency references contribute below $10^{-12}$.

\begin{figure}[htbp]
\centering
\includegraphics[width=0.45\textwidth]{figures/deviation_validation_with_full_data.png}
\caption{Validation of the empirical model for the refractive index of air at \SI{1762}{\nano\meter}. Main panel: Comparison of measured refractive index values with predictions from two models. Predictions from the empirical model developed in this work (red points) are compared against those from the published Mathar formulation \cite{mathar2007refractive} (black points, 1000-point subsample). The gold dashed line indicates the line of perfect agreement. The clustering of data points at higher and lower refractive index values reflects the seasonal distribution of the observations: a winter segment (January 2025, 8.3\% of the data) at lower temperatures and lower refractive index, and a summer segment (June--August 2025, 91.7\% of the data) at higher values, with no observations in the intervening months (see Supplemental Material). \textbf{(a)} Residual distribution for the empirical model: mean \num{-2.802e-14}$\pm$\num{1.837e-07} standard deviation ($R^2 = 0.996$). \textbf{(b)} Residual distribution for the Mathar model: mean \num{1.166e-06}$\pm$\num{1.930e-07} standard deviation ($R^2 = 0.829$); the offset is a static calibration offset removable by a single baseline calibration, after which the residual standard deviation is \num{1.930e-07}, comparable to that of the empirical model.}
\label{fig:validation}
\end{figure}

\section{Discussion and Conclusion}
\label{sec:discussion}

This work reports a quantum-referenced measurement of the refractive index of air at \SI{1762}{\nano\meter}, the $^{138}\mathrm{Ba}^+$ clock transition, over an eight-month campaign. The dataset of \num{145784} synchronized measurements yields the environmental coefficients $\partial n/\partial X$ for temperature, relative humidity, and pressure reported in Table~\ref{tab:coefficients}, together with their bootstrap standard errors. The residual analysis and the uncertainty budget of the empirical model relative to the Mathar formulation are reported in the Results section.

The comparison with the Mathar model is made on identical observed conditions: the model is evaluated at the same $(T,H,P)$ rows as the measurement and fitted with the same linear surrogate model, so that both fits share the same design matrix and no reference-temperature mismatch enters the comparison. Under this comparison the temperature and pressure coefficients agree with the Mathar surrogate to within \SI{0.35}{\percent} and \SI{0.93}{\percent}; the humidity coefficient is smaller in magnitude than the surrogate by \SI{24.8}{\percent}.

The humidity difference is statistically significant in the full campaign but not within the Mathar validity domain ($T \in [10,25]\,\si{\degreeCelsius}$). Three checks bear on its interpretation. First, the direction of the difference is opposite to the sign expected from anomalous dispersion near a water-vapor absorption line: such a response would increase the measured magnitude relative to the broadband value, whereas the measured magnitude is smaller. Second, a Kramers-Kronig evaluation of the HITRAN2020 water-vapor spectrum at the mean campaign conditions yields a coefficient of order $10^{-8}$\,\%RH$^{-1}$, without enhancement over the broadband model. Third, the magnitude of the difference is nearly spanned by the gain uncertainty of the BME280 humidity sensor: the \SI{4}{\percent} accuracy bound over the campaign temperature range corresponds to a worst-case gain deviation of 30.3\%, and a gain at this endpoint would account for $3.99 \times 10^{-9}$ of the $4.33 \times 10^{-9}$ difference, leaving $0.34 \times 10^{-9}$ (0.26 of the paired bootstrap standard error; see Supplemental Material). On the basis of these checks, the humidity difference is reported as an unresolved systematic, pending a measurement with an independently calibrated humidity sensor.

Several limitations of this work follow from the measurement configuration and are described in Supplemental Material. The environmental parameters were sampled by a single BME280 sensor at \SI{2}{\hertz}, so the residual reflects the spatial and temporal mismatch between that sensor and the \SI{4.2}{\meter} optical path, and turbulent fluctuations are not resolved. The humidity channel remains unresolved because no traceable humidity calibration was available; a gain systematic at the level of the sensor specification cannot be excluded. CO$_2$ variations were not tracked. These limitations do not affect the reported comparisons between the measured coefficients and the Mathar surrogate, which use identical environmental rows for both fits; they set the accuracy with which the coefficients themselves are determined, as quantified in Supplemental Material.

In summary, this work provides the following: a measurement architecture in which a single trapped ion sets the optical frequency of an interferometric refractometry channel over a \SI{4.2}{\meter} ambient air path; a multi-month, high-cadence dataset of \num{145784} synchronized measurements of refractive index and co-located environmental parameters; a comparison protocol that evaluates the Mathar model on the identical observed conditions and fits both the measurement and the model with the same linear surrogate; and an offset-fair assessment in which a single baseline calibration removes the static offset and the residual scatter of the empirical model agrees with the debiased Mathar model to within approximately \SI{5}{\percent}. The comparison also identifies the calibration of the environmental sensor as a limiting factor for the humidity channel.

\begin{acknowledgments}
We thank F. Hasse, H. Bao, P. Mullan and D. Palani for fruitful discussions, as well as A. Windt and I. Zhuravlev for comments on the manuscript. This research was funded by the European Research Council (ERC) under the European Union's Horizon 2020 research and innovation programme (grant number 648330), the Deutsche Forschungsgemeinschaft (DFG, grant number SCHA 973/9-1-3017959) and the Georg H. Endress Foundation. W.W. acknowledges financial support from the QUSTEC programme, funded by the European Union's Horizon 2020 research and innovation programme under the Marie Sklodowska-Curie (grant number 847471). T.S. acknowledges financial support from the DFG via the RTG DYNCAM 2717 and the Georg H. Endress foundation.
\end{acknowledgments}

\section*{Data Availability}
The data that support the findings of this study, including the synchronized time series of interferometric phase measurements and co-located environmental sensor readings, are openly available on Zenodo at Ref.~\cite{wu_2026_18800166}. Analysis scripts used to perform the multivariate regression, statistical diagnostics, and Kramers-Kronig integration are included in the same repository.

\bibliography{reference}

\end{document}